\subsection{Intrinsic minimal scalar winding}
\begin{statement}{P3K.38}{Proved}
Fix any of the thirty nonempty branches of P3F.53, including its actual
source coefficients, vertical line divisor, logarithmic norms, affine
cap maps, finite line characters, order-one clutches, and allowed
index-one fan. We write its line and source translation as
\[
 \mathcal M=\mathcal O(aG-O+D),\qquad Q=-bG,
 \qquad ab=d,
\]
with the notation and actual divisor comparisons of P3F.51--53. Let
\(m_i>1\) be its effective cap orders and put
\(M=\prod_i m_i\); an absent or order-one slot contributes one.
All norms \(n_i\) and the ordinary framing term are retained in
P3K.32's signed Seifert equation. In particular,
\(\gcd(m_i,n_i)=1\), and the nontrivial \(m_i\) are pairwise
coprime.

The actual scalar is the multiplier in
\[
                         B_\lambda=\lambda B_0.       \tag{IM.1}
\]
We continue the fixed construction on the logarithm cover of a small
punctured \(\lambda\)-disc. Its data at two logarithm values differing
by \(2\pi iN\) have a biholomorphic endpoint identification exactly
when
\[
                              N\in (M/b)\mathbb Z.     \tag{IM.2}
\]
This assertion allows every biholomorphic map between the resulting
threefolds. It imposes no preservation of a chosen elliptic marking,
line-circle action, source section, or toric chart on that map.
The positive integer \(M/b\) is therefore the least positive scalar
winding. Sufficiency is the actual family of P3K.35,
\[
                     \lambda=w^{M/b},\qquad w=r^b,
                     \lambda=r^M.                     \tag{IM.3}
\]
We prove necessity by intrinsic invariants of its multiple fibres.

For the component-supported OS49 family of P3K.34 and P3K.37, with
\(m=3k\), logarithmic vector
\(n(e_\gamma-e_u)+d_\delta\delta\), and \(n=\pm1\), the
corresponding least positive scalar winding is \(k\). Its actual
closed parameter is \(\lambda=w^k\), \(w=r^3\).
The argument for that statement uses its actual central Picard point
\(G'_0\); it uses no whole-fibre height formula.
\end{statement}
\begin{proof}
\par

\medskip\noindent\textbf{2. Every endpoint map preserves each nontrivial multiple fibre}

\par\nobreak

We first make the fibration intrinsic. Every threefold in question is
diffeomorphic to \(S^6\), by the integral topology of P3F.53 and
P3K.32, with the actual fan changes of P3F.35. Thus
\(H^2(Y;\mathbb R)=0\). A general fibre \(F\) of the displayed
map \(f:Y\to\mathbb P^1\) is a compact complex torus, hence is
Kähler. A horizontal effective divisor \(Z\) on \(Y\) would
restrict to a nonzero effective divisor \(Z|_F\) with zero real
first Chern class. For a Kähler form \(\omega\) on \(F\),
\[
 0=\int_F c_1(\mathcal O(Z)|_F)\wedge\omega
   =\int_{Z|_F}\omega>0,
\]
a contradiction. Every effective divisor is consequently vertical.

Let \(g\) be a meromorphic function on \(Y\). Its polar divisor
is vertical, so on a general compact torus fibre its restriction
has no polar divisor. Hartogs extension across its possible
codimension-two indeterminacy set makes that restriction holomorphic,
and hence constant. To see the descent at each base value, choose a
local coordinate \(q\). A sufficiently large power of \(q\circ f\)
cancels all divisorial poles on a small proper inverse image. The
resulting function is holomorphic by Hartogs. The fibres are connected;
the Stein factorization of this proper map to a normal disc therefore
gives \(f_*\mathcal O_Y=\mathcal O\). The holomorphic function
descends, and division by the same power of \(q\) gives a meromorphic
function on the disc. We have proved the literal field equality
\[
                    \mathbb C(Y)=f^*\mathbb C(t).       \tag{IM.4}
\]

A biholomorphism between endpoint threefolds induces an automorphism
of this field and hence a Möbius transformation of the base. The
equality of meromorphic maps extends over the entire threefold.
Divisorial multiplicities of its fibres are preserved. In the actual
cap chart the base function is \(q=z^{m_i}\), the root central
fibre is reduced, and the quotient is free with transitive component
action; its fibre is exactly \(m_i S_i\). The ordinary models
have divisorial gcd one. At the toric filling the same fact is the
primitivity of its physical valuation vector: the component orders
have gcd one, and changing the index-one fan retains that vector.
Thus the displayed effective caps are precisely the fibres of
multiplicity greater than one.

There are at most two such caps. Their coprime orders are distinct.
Every endpoint biholomorphism consequently carries each \(S_i\)
to the continuation of the same \(S_i\). This conclusion allows
permutations of ordinary fibres, including ones with the same Kodaira
type, and arbitrary source automorphisms. A biholomorphism of the
reduced supports lifts uniquely to their normalizations, by the
universal property of normalization. We may therefore use the
Albanese curve of a smooth cap support, or of its normalization,
without requiring any additional equivariance of the endpoint map.

\par

\medskip\noindent\textbf{3. The exact period lattice at a good-reduction cap}

\par\nobreak

Fix one good cap, and write its order and norm as \(m,n\).
The actual root line has central restriction
\[
                          N_0=\mathcal O_E(aG'_0-O).
                                                               \tag{IM.5}
\]
This is the line supplied by the root/quotient divisor diagrams in
P3F.41--42 and P3F.46--47. For every compatible vertical divisor
and fixed clutch, P3F.52's direct-image evaluation compares it with
(IM.5) by an actual base line and a fixed deck character. A base
line has trivial restriction to \(E\), and a character changes
the lift, not the underlying central line. Let
\[
                         t=\operatorname{ord}(aG'_0).
                                                               \tag{IM.6}
\]
This order includes the value one for the zero Picard point.

The literal minimal-root substitutions of P3F.51 give the following
points and automorphisms. At a good III or III* root, \(G'_0\)
is a nonzero point of order two and the deck derivative has order
four. At a good IV or IV* root it is a nonzero point of order three
and the derivative has order three. At a good II root the point
is \(O\), and the derivative has order six. In particular, with
\(o=\operatorname{ord}(G'_0)\),
\[
  t=\frac{o}{\gcd(o,a)},\qquad
  o=2,3,1\quad\hbox{in those three cases}.             \tag{IM.7}
\]
The program retains this multiplication by \(a\). For example the
good E7 line with \(a=2\) has order one, whereas that with \(a=1\)
has order two.

Write the central source cap map as \(z\mapsto Dz+R\) on
\(E=\mathbb C/\Lambda_E\). Here \(D\ne1\), its order divides
\(m\), and \(1+D+\cdots+D^{m-1}=0\). The source part of
\(A^m=B^n\) gives \(nbG'_0=0\). We also have \(mG'_0=0\).
Since \(\gcd(m,n)=1\), the Bezout identity gives \(bG'_0=0\).
The actual \(B_0\) therefore covers the identity on \(E\), and
acts on \(N_0\) by a single nonzero constant \(c\). Put
\[
                  \tau=\frac{\log(c\lambda)}{2\pi i}.
                                                               \tag{IM.8}
\]

We now exhibit every period rather than just an invariant subspace.
The torsion line (IM.5) has its finite unitary character
\(\chi:\Lambda_E\to\mu_t\), of exact image order \(t\).
In a logarithmic line coordinate \(\eta\), its universal cover is
\(\mathbb C^2\), with generators
\[
 (z,\eta)\longmapsto(z+\omega,\eta+\theta_\omega),
 \quad e^{2\pi i\theta_\omega}=\chi(\omega),
 \qquad (z,\eta)\longmapsto(z,\eta+1).
                                                               \tag{IM.9}
\]
Here every \(\theta_\omega\) is real, and their classes together
with \(\mathbb Z\) generate exactly \((1/t)\mathbb Z\).
The quotient by \(B_\lambda\) adds the period \((0,\tau)\)
and is the actual compact central abelian surface.

The equality \(D^*N_0=N_0\) makes the finite unitary character
\(D\)-invariant. Thus an affine source map \(Dz+R\) has a
constant line lift in (IM.9). Every other holomorphic lift differs
by an invertible holomorphic function on the compact elliptic curve,
which is constant. The actual cap generator can consequently be
written, with its fixed character retained, as
\[
 A_\lambda(z,\eta)
       =(Dz+R,\eta+n\tau/m+\kappa).                  \tag{IM.10}
\]
The exact norm gives \(m\kappa\in\mathbb Z\): the source lift
has zero translation after \(m\) iterations because the geometric
sum of the powers of \(D\) is zero, and a deck period with zero
source coordinate has integral line coordinate. We retain \(\kappa\) in every
formula below; no finite character is discarded.

All holomorphic one-forms on the abelian cover are constant.
The affine linear part \(\operatorname{diag}(D,1)\) fixes just
the line generated by \(d\eta\). Integration over the full
affine deck group (IM.9)--(IM.10) gives its complete period image
\[
 \mathcal P=(1/t)\mathbb Z+\tau\mathbb Z
                         +(n\tau/m+\kappa)\mathbb Z. \tag{IM.11}
\]
Choose integers \(x,y\) with \(mx+ny=1\). Then the actual
period combination
\(x\tau+y(n\tau/m+\kappa)\) is
\(\beta=\tau/m+y\kappa\), and conversely
\[
  \tau=m\beta-my\kappa,\qquad
  n\tau/m+\kappa=n\beta+mx\kappa.
                                                               \tag{IM.12}
\]
The remainders belong to \((1/t)\mathbb Z\). These equations
prove the lattice equality, with its arrows in both directions,
\[
                 \mathcal P=(1/t)\mathbb Z
                             +(\tau/m+y\kappa)\mathbb Z.
                                                               \tag{IM.13}
\]
Replacing \((x,y)\) by \((x+ns,y-ms)\) changes the second
generator by the first lattice, so the result is independent of
the Bezout representative.

The map induced by \(\eta\) to \(\mathbb C/\mathcal P\)
is the Albanese map. This follows directly from the invariant
one-form and its complete integration lattice; equivalently its
fibre is the elliptic cover with period kernel \(\ker\chi\).
Multiplication by \(t\) gives its Tate coordinate
\[
           Q_i=e^{2\pi i t y\kappa}(c\lambda)^{t/m}.
                                                               \tag{IM.14}
\]
The norm \(n\), source affine term, and fixed line character occur
in the displayed constant. The exponent is determined by the actual
Picard point (IM.6).

\par

\medskip\noindent\textbf{4. The normalization and its period at a reduced-root cap}

\par\nobreak

We now use an actual reduced root of type II, III, IV, or I2. Its
central curve has, respectively, \(r=1,2,3,2\) rational components
in its normalization. We retain the precise root line, whose degree
on each component is zero by the line diagrams in P3F.36--49 and
P3F.52. Its pullback to each \(\mathbb P^1\) is therefore
trivial. This assertion does not trivialize the line on the singular
curve, nor discard its nodal or infinitesimal gluing data.

On the actual root line the two maps satisfy
\[
 B_\lambda=\lambda B_0,\qquad
 A_\lambda=\lambda^{n/m}A_0,\qquad
 A_\lambda^m=B_\lambda^n,
 \qquad[A_\lambda,B_\lambda]=1.                       \tag{IM.15}
\]
The fractional power here is on the chosen logarithm cover of the
scalar disc. For integers \(x,y\) with \(mx+ny=1\), form the
literal holomorphic line map
\[
                         H=B_\lambda^xA_\lambda^y.
                                                               \tag{IM.16}
\]
Negative powers are allowed because both maps are invertible in
this cap. Their exact norm and commutativity prove
\[
 H^m=B_\lambda^{mx+ny}=B_\lambda,\qquad
 H^n A_\lambda^{-1}
        =(B_\lambda^nA_\lambda^{-m})^x=1.
                                                               \tag{IM.17}
\]
Consequently their actual action group is precisely \(\langle H\rangle\).

The sphere cap's component action is transitive, as established by
the integral cap constructions. Since it is now generated by \(H\),
that generator cycles all \(r\) normalization components. Choose
fixed trivializations of their degree-zero lines and label the
components in this cyclic order. On them the literal maps have form
\[
 (z,V,j)\longmapsto
        (h_j(z),c_j\lambda^{1/m}V,j+1).               \tag{IM.18}
\]
The source maps \(h_j\) are Möbius transformations, independent
of \(\lambda\). The coefficients \(c_j\) are constant in
\(z\): an invertible map between the two trivial lines on compact
\(\mathbb P^1\) is a nowhere-zero holomorphic function on that
curve. The exponent is exactly
\(x+yn/m=1/m\). The constants retain every finite character and
the fixed singular-line gluing.

The stabilizer of one component is \(\langle H^r\rangle\).
Its quotient is therefore the explicit surface
\[
  \widetilde S_i=
  (\mathbb P^1\times\mathbb C^*)/\langle
       (z,V)\longmapsto(h(z),c_i\lambda^{r/m}V)
                                      \rangle .       \tag{IM.19}
\]
This is the normalization of the reduced cap support. To check the
map, normalize the root central curve and pull back its line. The
resulting finite map is equivariant for (IM.16). For sufficiently
small \(|\lambda|\), \(H^r\) strictly contracts the line norm,
so its action is free and properly discontinuous. Quotienting the
finite equivariant map, using compact fundamental annuli, gives a
finite map from (IM.19) to the actual central support. It is
birational because it is an isomorphism over the smooth locus of
the root central curve. The source is normal, so the finite map is
its normalization. This establishes the quotient arrow, including
the cuspidal type II and the multiple-intersection type III/IV
central curves.

The actual projection \([(z,V)]\mapsto[V]\) makes (IM.19) a
holomorphic \(\mathbb P^1\)-bundle over
\[
            E_{Q_i}=\mathbb C^*/Q_i^{\mathbb Z},
            \qquad Q_i=c_i\lambda^{r/m}.             \tag{IM.20}
\]
All holomorphic one-forms come from this elliptic curve, because
they vanish along the \(\mathbb P^1\) fibres. The homotopy exact
sequence of the displayed locally trivial bundle gives an
isomorphism on fundamental groups. Thus its full integration
lattice, and its Albanese curve, are exactly those of (IM.20).
In particular the source Möbius monodromy \(h\) introduces no
unrecorded isogeny into the Albanese period.

For the quadratic I1* cap, \(r=m=2\). Here \(B\) preserves
the two components while \(A\) exchanges them; (IM.17) still
cycles them and gives \(Q_i=c_i\lambda\). This cap imposes
no scalar-cover factor on its own. The reduced III and IV root
progressions have \(\gcd(m,r)=1\); they retain the full factor
\(m\). The type II normalization has \(r=1\).

\par

\medskip\noindent\textbf{5. Intrinsic obstruction to every smaller winding}

\par\nobreak

For a small nonzero Tate coordinate we have
\(j(Q)=Q^{-1}+744+O(Q)\), with the same normalization as in
P3K.36. For any root of unity \(\zeta\ne1\),
\[
      Q\bigl(j(Q)-j(\zeta Q)\bigr)
                         \longrightarrow1-\zeta^{-1}\ne0.
                                                               \tag{IM.21}
\]
There are only finitely many phase ratios for each fixed cap and
finitely many caps. We shrink the scalar disc once so that all the
nontrivial differences in (IM.21) are nonzero there.

After scalar winding \(N\), formula (IM.14) multiplies \(Q_i\)
by \(e^{2\pi iNt_i/m_i}\) at a good cap, and formula (IM.20)
multiplies it by \(e^{2\pi iNr_i/m_i}\) at a reduced cap. An
endpoint biholomorphism would give an isomorphism of the intrinsic
Albanese curves in Section 2, so it would make their \(j\)-invariants
equal. Equations (IM.21) therefore force, respectively,
\[
 m_i\mid Nt_i,
 \qquad m_i\mid Nr_i.                                \tag{IM.22}
\]
Thus the exact local necessary denominators are
\[
 d_i=\frac{m_i}{\gcd(m_i,t_i)}\quad\hbox{for a good root},
 \qquad
 d_i=\frac{m_i}{\gcd(m_i,r_i)}\quad\hbox{for a reduced root}.
                                                               \tag{IM.23}
\]
Each \(d_i\) divides \(m_i\), so distinct \(d_i\) are coprime.
The combined necessary winding is their product.

The actual root cases have the following removed factors:

\begin{longtable}[]{@{}llr@{}}
\toprule\noalign{}
Root type & Actual central datum & Factor \(m_i/d_i\) \\
\midrule\noalign{}
\endhead
\bottomrule\noalign{}
\endlastfoot
Good III or III* & \(t_i=2/\gcd(2,a)\) & \(t_i\) \\
Good IV or IV* & \(t_i=3/\gcd(3,a)\) & \(t_i\) \\
Good II & \(t_i=1\) & 1 \\
Reduced III, orders \(4j+1\) or \(4j+3\) & \(r_i=2\) & 1 \\
Reduced IV, orders \(3j+1\) or \(3j+2\) & \(r_i=3\) & 1 \\
Reduced II & \(r_i=1\) & 1 \\
Quadratic I1* root I2 & \(m_i=r_i=2\) & 2 \\
Absent or order-one slot & ordinary marked cap & 1 \\
\end{longtable}

In every nonempty F53 branch the product of these factors is exactly
the specified \(b\). The complete branch list and the exact
residue-class verification are recorded below and in the certificate.
Consequently
\[
                       \prod_i d_i=M/b.                \tag{IM.24}
\]
This proves the necessity in (IM.2). The quotient constructed in
P3K.35 proves its sufficiency, including every retained order-one
clutch and all allowed fans. Pulling that actual family back by
\(w\mapsto w^q\), or reversing the continuation direction,
gives the endpoint identifications for every integer multiple.
No statement about the mapping class of an additional endpoint
automorphism is required or deduced.

\par

\medskip\noindent\textbf{6. The component-supported OS49 cap}

\par\nobreak

The good IV* root of P3F.26 has the actual central line
\(\mathcal O(G'_0-O)\), where \(G'_0\) is a nonzero
three-torsion point. Its order is \(m=3k\). The component-supported
line at the I2 fibre changes neither this Picard point nor the
central IV* root diagram. The cap calculation of Section 3 therefore
has \(t=3\). Its intrinsic Albanese Tate parameter has exponent
\(3/m=1/k\).

For completeness we retain the mixed logarithmic marking directly.
With the actual covector marking of P3F.26, invariant covectors are
\(\psi\) and \(h=-w+3\delta\). In the coherent central period
matrix
\[
 \Pi=[(-3\mu,\beta),(-\tau_E,-\mu),(1,0),(0,1)],
 \qquad \tau_E^2-3\tau_E+3=0,\quad \mu=\tau_E/3,
\]
the invariant holomorphic form is \(-da+3db\). Its period row
is \(h+T\psi\), where
\[
 T=3\beta+\tau_E,
 \qquad \beta=\frac{\log\lambda}{2\pi i}+c_0.
\]
The actual affine vector
\(n(e_\gamma-e_u)+d_\delta\delta\) adds
\((nT+3d_\delta)/m\). For \(n=\pm1\), the complete period
lattice is
\[
 \mathbb Z+\mathbb Z\frac{T+3nd_\delta}{m}.
                                                               \tag{IM.25}
\]
Indeed the added period is \(n\) times the displayed second
generator, and \(T\) differs from \(m\) times that generator
by the integer \(3nd_\delta\). The generator \(h\) gives the
primitive real period one. This proves directly that the fixed
delta term and either norm sign affect the constant, while
\[
                        Q_i=c_i\lambda^{1/k}.          \tag{IM.26}
\]
The algebraic-reduction argument of Section 2 applies to these
actual spheres as well. Their unique nontrivial multiple fibre has
order \(3k\), so every endpoint map preserves it. Equation (IM.21)
forces \(k\mid N\). P3K.37 supplies the exact quotient
\(\lambda=w^k\), proving sharpness. This conclusion uses the
actual cap periods; the component line square remains the separate
value \(-2k\) established in P3K.34.
\end{proof}
\par

\medskip\noindent\textbf{7. Executed certificate and complete branch coverage}

\par\nobreak

The program \path{code/part3/scalar_cap_minimality.py} verifies the
two-way good-cap period basis, the exact cyclic-generator identities,
the multiplication of the actual point by \(a\), and all gcd
formulas on every residue class modulo the finite central order.
It checks all thirty nonempty branches against \(M/b\), and
retains the impossible good III*/good II branch as excluded by its
common even cap orders. It does not replace a geometric period map
by a numerical inference. Its 171 exact checks pass; full output is
in the matching JSON and OUT files.

The parameters in the following table retain all positivity bounds,
coprimality conditions, norm equations, and finite characters in the
corresponding F53 branch. Its two denominators are attached to the
same two actual cap slots; a one denotes an absent or unconstraining
slot. The table is generated directly from the audited F53 catalogue.

\begin{longtable}[]{@{}
  >{\raggedright\arraybackslash}p{(\columnwidth - 8\tabcolsep) * \real{0.1765}}
  >{\raggedright\arraybackslash}p{(\columnwidth - 8\tabcolsep) * \real{0.1765}}
  >{\raggedleft\arraybackslash}p{(\columnwidth - 8\tabcolsep) * \real{0.2353}}
  >{\raggedright\arraybackslash}p{(\columnwidth - 8\tabcolsep) * \real{0.1765}}
  >{\raggedleft\arraybackslash}p{(\columnwidth - 8\tabcolsep) * \real{0.2353}}@{}}
\toprule\noalign{}
\begin{minipage}[b]{\linewidth}\raggedright
Branch
\end{minipage} & \begin{minipage}[b]{\linewidth}\raggedright
Source profile
\end{minipage} & \begin{minipage}[b]{\linewidth}\raggedleft
(a,b)
\end{minipage} & \begin{minipage}[b]{\linewidth}\raggedright
Cap denominators
\end{minipage} & \begin{minipage}[b]{\linewidth}\raggedleft
Minimum winding
\end{minipage} \\
\midrule\noalign{}
\endhead
\bottomrule\noalign{}
\endlastfoot
CS01 & \(III^*+3I_{1}\) & (1,2) & \(2 l\), \(1\) & \(2 l\) \\
CS02 & \(III^*+3I_{1}\) & (2,1) & \(4 l + 3\), \(1\) & \(4 l + 3\) \\
CS03 & \(III^*+II+I_{1}\) & (1,2) & \(2 l\), \(6 k + 1\) & \(2 l \left(6 k + 1\right)\) \\
CS05 & \(III^*+II+I_{1}\) & (2,1) & \(4 l + 3\), \(6 k + 1\) & \(\left(6 k + 1\right) \left(4 l + 3\right)\) \\
CS06 & \(III^*+II+I_{1}\) & (2,1) & \(4 l + 3\), \(6 k\) & \(6 k \left(4 l + 3\right)\) \\
CS07 & \(I_{8}+4I_{1}\) & (8,1) & \(1\), \(1\) & \(1\) \\
CS08 & \(I_{8}+II+2I_{1}\) & (8,1) & \(6 l + 1\), \(1\) & \(6 l + 1\) \\
CS09 & \(I_{8}+II+2I_{1}\) & (8,1) & \(6 l\), \(1\) & \(6 l\) \\
CS10 & \(I_{7}+I_{2}+3I_{1}\) & (14,1) & \(1\), \(1\) & \(1\) \\
CS11 & \(I_{7}+I_{2}+II+I_{1}\) & (14,1) & \(6 l + 1\), \(1\) & \(6 l + 1\) \\
CS12 & \(I_{7}+I_{2}+II+I_{1}\) & (14,1) & \(6 l\), \(1\) & \(6 l\) \\
CS13 & \(I_{7}+III+2I_{1}\) & (7,2) & \(2 l\), \(1\) & \(2 l\) \\
CS14 & \(I_{7}+III+2I_{1}\) & (14,1) & \(4 l + 1\), \(1\) & \(4 l + 1\) \\
CS15 & \(IV^*+I_{2}+2I_{1}\) & (2,3) & \(l\), \(1\) & \(l\) \\
CS16 & \(IV^*+I_{2}+2I_{1}\) & (6,1) & \(3 l + 2\), \(1\) & \(3 l + 2\) \\
CS17 & \(IV^*+III+I_{1}\) & (1,6) & \(l\), \(2 k\) & \(2 k l\) \\
CS18 & \(IV^*+III+I_{1}\) & (2,3) & \(l\), \(4 k + 1\) & \(l \left(4 k + 1\right)\) \\
CS19 & \(IV^*+III+I_{1}\) & (3,2) & \(3 l + 2\), \(2 k\) & \(2 k \left(3 l + 2\right)\) \\
CS20 & \(IV^*+III+I_{1}\) & (6,1) & \(3 l + 2\), \(4 k + 1\) & \(\left(4 k + 1\right) \left(3 l + 2\right)\) \\
CS21 & \(I_{1}^*+I_{3}+2I_{1}\) & (6,2) & \(1\), \(1\) & \(1\) \\
CS22 & \(I_{1}^*+IV+I_{1}\) & (2,6) & \(1\), \(k\) & \(k\) \\
CS23 & \(I_{1}^*+IV+I_{1}\) & (6,2) & \(1\), \(6 l + 1\) & \(6 l + 1\) \\
CS24 & \(I_{5}+I_{4}+3I_{1}\) & (20,1) & \(1\), \(1\) & \(1\) \\
CS25 & \(I_{5}+I_{4}+II+I_{1}\) & (20,1) & \(6 l + 1\), \(1\) & \(6 l + 1\) \\
CS26 & \(I_{5}+I_{4}+II+I_{1}\) & (20,1) & \(6 l\), \(1\) & \(6 l\) \\
CS27 & \(I_{5}+I_{3}+I_{2}+2I_{1}\) & (30,1) & \(1\), \(1\) & \(1\) \\
CS28 & \(I_{5}+I_{3}+III+I_{1}\) & (15,2) & \(2 l\), \(1\) & \(2 l\) \\
CS29 & \(I_{5}+I_{3}+III+I_{1}\) & (30,1) & \(4 l + 1\), \(1\) & \(4 l + 1\) \\
CS30 & \(I_{5}+IV+I_{2}+I_{1}\) & (10,3) & \(l\), \(1\) & \(l\) \\
CS31 & \(I_{5}+IV+I_{2}+I_{1}\) & (30,1) & \(3 l + 1\), \(1\) & \(3 l + 1\) \\
\end{longtable}
